\documentclass[12pt]{article}
\usepackage[T1]{fontenc}
\usepackage[utf8]{inputenc}
\usepackage{lmodern}
\usepackage{textcomp}
\usepackage{enumitem}
\usepackage{geometry}
\usepackage{graphicx}
\usepackage{amsmath, amssymb}
\geometry{margin=1in}
\begin{document}
\begin{center}
Astronomy - Graduate Level Assessment 15 \
Total Marks: 40 \
Answer all questions.
\end{center}

\section*{Multiple Choice (10 marks)}
\begin{enumerate}[label=\arabic*.]
\item Suppose you live on the Moon. How long is a day (i.e. from sunrise to sunrise)?
\begin{enumerate}[label=(\alph*)]
    \item about 18 years
    \item 24 hours
    \item 29 Earth days
    \item a year
\end{enumerate}
\item 20000 years from now ...
\begin{enumerate}[label=(\alph*)]
    \item The Moon will be closer to the Earth and the Earth's day will be longer.
    \item The Moon will be closer to the Earth and the Earth's day will be shorter.
    \item The Moon will be further from the Earth and the Earth's day will be longer.
    \item The Moon will be further from the Earth and the Earth's day will be shorter.
\end{enumerate}
\item Approximately how old is the surface of Venus?
\begin{enumerate}[label=(\alph*)]
    \item 750 million years.
    \item 2 billion years.
    \item 3 billion years.
    \item 4.5 billion years.
\end{enumerate}
\item Which of the following is/are true?
\begin{enumerate}[label=(\alph*)]
    \item Titan is the only outer solar system moon with a thick atmosphere
    \item Titan is the only outer solar system moon with evidence for recent geologic activity
    \item Titan's atmosphere is composed mostly of hydrocarbons
    \item A and D
\end{enumerate}
\item You've made a scientific theory that there is an attractive force between all objects. When will your theory be proven to be correct?
\begin{enumerate}[label=(\alph*)]
    \item The first time you drop a bowling ball and it falls to the ground proving your hypothesis.
    \item After you've repeated your experiment many times.
    \item You can never prove your theory to be correct only "yet to be proven wrong".
    \item When you and many others have tested the hypothesis.
\end{enumerate}
\end{enumerate}

\section*{True / False (10 marks)}
\textit{Answer True or False}
\begin{enumerate}[label=\arabic*.]
\setcounter{enumi}{5}
\item True or False: Barnard's Star is the second closest star system to Earth after Proxima Centauri.
\item True or False: Water-ice is uniformly distributed across Mercury and is not limited to the polar regions.
\item True or False: For a total solar eclipse to be observed, the Moon's umbra must intersect the observer's location on Earth.
\item True or False: The Schwarzschild radius is primarily used to describe the stellar structure of red dwarfs.
\item True or False: Pluto takes approximately 250 years to complete one orbit around the Sun.
\end{enumerate}

\section*{Long Form Response (20 marks)}
\textit{Answer each question with a clear and complete explanation.}
\begin{enumerate}[label=\arabic*.]
\setcounter{enumi}{10}
\item Discuss the implications of the Moon's orbital plane aligning with the ecliptic plane, particularly focusing on the frequency and nature of solar eclipses.
\item Discuss the significance of temperature measurement in astronomy, specifically explaining the conversion from Celsius to Kelvin, and analyze why Kelvin is preferred for astronomical contexts.
\end{enumerate}

\vfill
\noindent\textit{End of Paper}
\end{document}